\chapter{Governance Policies}
\label{chap:Governance_Policies}

\section{Purpose}

Governance Policies are the concrete rules that regulate:
\begin{itemize}
    \item which Engrams are allowed or forbidden,
    \item where and how the CME may be deployed,
    \item who may perform which operations,
    \item how HITL mechanisms must behave in normal and emergency modes.
\end{itemize}

They translate P-class charters and EAC decisions into operational constraints.

\section{Policy Taxonomy}

We distinguish:

\subsection{Access Policies}

Regulate:
\begin{itemize}
    \item who may interact with the CME,
    \item which roles can invoke which SIGIL Pathways,
    \item which contexts (domains, networks, jurisdictions) are allowed.
\end{itemize}

\subsection{Engram Policies}

Regulate:
\begin{itemize}
    \item Engram types that may be stored in GEL,
    \item forbidden Engram categories (e.g., certain weapons, hate content),
    \item conditions under which Engrams may be retired or quarantined.
\end{itemize}

\subsection{Deployment Policies}

Regulate:
\begin{itemize}
    \item which environments may host the CME,
    \item data boundary and privacy requirements,
    \item resource limits and safety thresholds.
\end{itemize}

\subsection{Incident Policies}

Regulate:
\begin{itemize}
    \item detection and response to misalignment or incidents,
    \item emergency halt and rollback procedures,
    \item reporting obligations and timelines.
\end{itemize}

\section{Policy Schema}

A policy $P$ is defined as:
\[
P = \langle id,\, type,\, scope,\, rule,\, enforcement,\, review\_cycle \rangle
\]

Where:
\begin{itemize}
    \item $id$: unique identifier,
    \item type: access / Engram / deployment / incident / other,
    \item scope: who/what/where the policy applies to,
    \item rule: natural-language description,
    \item enforcement: how the harness enforces it,
    \item review\_cycle: how often it must be re-evaluated.
\end{itemize}

\section{Enforcement Mechanisms}

Governance Policies are enforced through:

\begin{itemize}
    \item SIGIL Pathways (for high-impact operations),
    \item automatic harness checks (for routine operations),
    \item CAR review (for ambiguous or high-risk cases),
    \item EAC decisions (for moral/ethical conflicts).
\end{itemize}

Violations trigger:
\begin{itemize}
    \item alarms in the Operator Interface,
    \item entries in the HITL audit log,
    \item potential incident escalation.
\end{itemize}

\section{Change Management}

Policy changes follow:

\begin{enumerate}
    \item Proposal (from operators, stakeholders, or CME telemetry).
    \item CAR classification (Class B or A).
    \item Review and deliberation (CAR and possibly EAC).
    \item Implementation via:
    \begin{itemize}
        \item updated policy records,
        \item new or updated SIGIL Pathways,
        \item harness configuration changes.
    \end{itemize}
    \item Logging and future review scheduling.
\end{enumerate}

\section{Summary}

Governance Policies operationalize high-level values and constraints into
concrete, enforceable rules that guide CME behaviour across all environments
and time scales.
